% Cover, references, and appendix are excluded from the four-page body.
\documentclass[12pt,a4paper]{article}
\usepackage[margin=2.5cm]{geometry}
\usepackage{fontspec}
\setmainfont{Times New Roman}
\usepackage{array,tabularx,booktabs}
\usepackage{enumitem}
\usepackage{titlesec}
\usepackage{fancyhdr}
\usepackage[hidelinks]{hyperref}
\urlstyle{rm}
\DeclareUrlCommand\path{\urlstyle{rm}}
\usepackage{microtype}
\setlength{\parindent}{0pt}
\setlength{\parskip}{5pt}
\setlength{\emergencystretch}{2em}
\setlist[itemize]{leftmargin=*,nosep,topsep=3pt}
\renewcommand{\arraystretch}{1.12}
\newcolumntype{Y}{>{\raggedright\arraybackslash}X}
\titleformat{\section}{\large\bfseries}{\thesection}{0.6em}{}
\titleformat{\subsection}{\normalsize\bfseries}{\thesubsection}{0.6em}{}
\titlespacing*{\section}{0pt}{8pt}{5pt}
\titlespacing*{\subsection}{0pt}{6pt}{3pt}
\pagestyle{fancy}
\fancyhf{}
\fancyhead[L]{COMP5424 -- Group 7}
\fancyhead[R]{Phase 2: Prototype}
\fancyfoot[C]{\thepage}
\setlength{\headheight}{15pt}
\newcommand{\code}[1]{\textnormal{#1}}
\newcommand{\repo}{https://github.com/itemgroup/COMP5424_CampusEmergency_XR}

\begin{document}
\begin{titlepage}
\centering
\vspace*{1.1cm}
{\Large COMP5424 Extended Reality\par}
\vspace{0.8cm}
{\LARGE\bfseries Prototype Design and Technical Note\par}
\vspace{0.5cm}
{\Large Phase 2 -- Group 7\par}
\vspace{0.8cm}
{\large Scenario D: Campus Emergency Preparedness\par}
\vspace{0.3cm}
{\large VR Stair-Evacuation and Lift-Safety Practice\par}
\vspace{0.8cm}
\begin{tabularx}{\textwidth}{@{}lY@{}}
\textbf{Virtual client} & Harbourview University Safety and Resilience Office \\
\textbf{Submission date} & 17 October 2026 \\
\textbf{Repository} & \url{\repo} \\
\textbf{Source baseline} & \code{final-submission-2026-10-17} at \code{9f2b3ff} \\
\textbf{Submitted build} & \path{Build/CampusEmergency.zip}; Windows x86\_64 IL2CPP, source \code{9f2b3ff} \\
\textbf{Target-PC evidence} & Team-reported Windows PCVR test on PQ604B / 604B205; Quest Link/Air Link; 17 October 2026 \\
\end{tabularx}
\vspace{0.8cm}
\begin{tabularx}{0.75\textwidth}{@{}YY@{}}
\toprule
\textbf{Member} & \textbf{Student ID} \\
\midrule
SUN Jiashi & 26044457G \\
SONG Zheyi & 26064683G \\
WENG Zhirui & 26101594G \\
FENG Kedi & 26058499G \\
CHU Silu & 26115249G \\
ZHU Yu & 26101969G \\
\bottomrule
\end{tabularx}
\vfill
\textit{A supervised educational demonstration; not official emergency guidance,
an approved campus evacuation plan, or certification.}
\end{titlepage}
\pagenumbering{arabic}

% BODY PAGE 1
\section{Prototype Scope and Requirements Traceability}
\subsection{Purpose and implemented scope}
Phase 1 proposed sign interpretation and recovery when an initial evacuation
route becomes unavailable~\cite{phase1}. Phase 2 prioritises a runnable,
testable XR journey rather than visual polish~\cite{brief}. The current source
implements \textbf{stair-evacuation and lift-safety practice} with an optional,
authored route-closure and recovery exercise.

\textbf{Implemented journey:} Enter System $\rightarrow$ welcome and task
briefing $\rightarrow$ stair route and fire alert $\rightarrow$ marked
landings $\rightarrow$ 2F practice endpoint and persistent completion card.
Optional reroute mode first presents the original route, then closes it at the
3F decision corner and requires a signposted same-floor alternative before
rejoining the stairs. The diagram and controller tutorial remain optional.
A separate G-floor courtyard assembly interaction is additional content, not
the 2F route's completion condition.

\subsection{Requirement-to-feature mapping}
Status describes software scope, not demonstrated learning effectiveness.
Evidence keys refer to Appendix A and the source baseline~\cite{implementation}.

\begin{tabularx}{\textwidth}{@{}p{0.65cm}YY@{}}
\toprule
\textbf{ID} & \textbf{Requirement / implementation} & \textbf{Status / evidence} \\
\midrule
R1 & Complete evacuation practice journey: briefing, stairs, endpoint and optional recovery. &
Implemented authored practice; learning effect not validated. S1--S3; V1, V3. \\
R2 & Interpret signs and environmental cues: floor/lift signs, diagram and task card. &
Implemented cues; independent interpretation not established. S1--S3; V2. \\
R3 & Recover when an initial route becomes unavailable. &
Optional 3F closure, enforced alternative and stair rejoin; authored, not dynamic routing. V3. \\
R4 & Explain why an action or route is appropriate. &
Lift warning and safety review explain stair/lift choice and closure cues; learning effect unmeasured. S2--S4; V3. \\
R5 & Calm, inclusive, low-risk practice. &
Partial: text feedback, optional tutorial, teleport and moderate audio settings; comfort needs user evidence. V4. \\
R6 & Pause, exit and restart safely. &
Implemented; PCVR pass reported for PQ604B/604B205; broader device coverage remains unverified. S2; V5, V6. \\
\bottomrule
\end{tabularx}

\subsection{Changes from Phase 1}
The route-change requirement is implemented as a bounded, optional exercise
using one fixed closure, signs and an authored corridor, not as dynamic
routing. Reusing the VR template and imported building assets enabled
integration of movement, signs and feedback. The environment includes multiple
floors, exceeding the proposed small scene. Campus geometry is illustrative,
not a validated evacuation map. The prototype can be evaluated
for task completion and cue usability, but current formative evidence cannot
establish learning or transfer. AI co-navigation, multiplayer, certification
and graphic emergency effects remain outside scope.

\newpage
% BODY PAGE 2
\section{Interaction and Prototype Design}
\subsection{Core interaction and route decisions}
The welcome NPC introduces the demonstration. Continue reveals the task list,
evacuation diagram and optional five-step controller practice. Users then read
corridor signs and enter the marked stair route. Arrival at its starting
waypoint starts the exercise audio; ordered checkpoints update the route task
state. Position-based triggers support both walking and teleport arrival.

Four floors contain eight lift signs. Nearby panels appear within 2.5\,m
horizontally and 1.25\,m vertically of a sign. During the fire exercise,
entering this area can announce that lifts must not be used; activating a panel
adds warning text and an error cue. Panels sit 0.5\,m above and 0.5\,m along
world $+Z$ from each sign to reduce surface intersection.

\textbf{Decision and rationale:} Fixed markers make arrival feedback repeatable
and avoid dependence on NPC placement. The optional reroute exercise closes
the initial 3F route, blocks invalid teleport destinations and cross-barrier
shortcuts, then signposts an opposite same-floor corridor back to the existing
stairs. It is an authored scenario, not dynamic route planning. Outside this
exercise ordinary exploration is unchanged. The diagram's arrows and automatic
lift announcement reveal directions, so they confound claims of unaided choice.

\subsection{Locomotion, controls, and feedback}
\textbf{Input:} The left stick moves; the right stick turns and its upward
teleport input aims at a valid target. Release confirms teleport, the right
grip cancels, and the right trigger interacts with supported objects/UI.
Teleport reduces reliance on continuous travel; individual comfort has not
been evaluated. The right stick uses 45-degree snap turning.

\textbf{UI:} Tutorial/navigation panels are world-fixed after opening.
Nearer/Farther physically move a panel in 0.2\,m steps (normal range
0.65--2.2\,m); Review reopens the guide at the current viewpoint. Thus the
panel does not chase head movement, although users can leave it behind.
The completion card persists after arrival at the endpoint or End NPC,
independently of Continue. It offers a safety review, a Continue button that
leaves the current state unchanged, and a Back action that reloads the scene
to clear the previous attempt and open the classroom welcome page. The review
describes stairs, lift avoidance, noticing a closure and the simulated
endpoint; it is not evidence of retained learning.

\textbf{Pause/reset:} Pause sets the simulation time scale to zero and opens
a view-following menu; Resume restores it. Restart reloads the active scene;
Exit quits the player. Input disabling while paused has not been separately
verified.

\subsection{Calmness, accessibility, and journey evidence}
English text, symbols, high-contrast panels and adjustable panel distance
provide visual support alongside audio. Low-intensity cues avoid deliberate
flashes, smoke or graphic harm. Multilingual access, seated reach, silent-mode
comprehension and individual comfort are not yet established.

\newpage
% BODY PAGE 3
\section{Technical Implementation and Functional Verification}
\subsection{Architecture and major technical decisions}
\textbf{Configuration:} Unity 6000.3.25f1; URP 17.3.0; XR Interaction Toolkit
3.5.1; Input System 1.20.0; OpenXR 1.17.1; Android IL2CPP/ARM64 for
Quest 3~\cite{implementation}; the Windows PCVR test is summarised below.

\textbf{Control:} ControllerGuidePanel tracks route start, stair entry,
landings and completion from teleport-anchor positions; a separate 2.5\,m
check shows the End NPC card once without advancing tasks. Opt-in
RouteChangeExercise adds the selection filter, closure barrier and alternative
signs. NpcGuideDialogue manages dialogue cards; MainMenuController restores
campus roots and reloads the classroom; CampusAudioDirector and LiftSafetyAudio
manage audio.

\textbf{Trade-offs:} One scene preserves the XR rig; authored markers and
offline speech reduce integration and network risk. Recovery is limited to
one optional authored closure and corridor. The ordinary endpoint card does
not require ordered completion; reroute practice does.

\subsection{Executable and repository reproducibility}
The Windows Standalone x86\_64 IL2CPP build from \code{9f2b3ff} is in
\path{Build/CampusEmergency.zip}. The team reported launching it on PQ604B
machine 604B205 on 17 October 2026 via Quest Link/Air Link, with Meta Quest
Link as the active OpenXR Runtime; all intended features reportedly worked.
No task-by-task result log was retained. This records one PCVR setup, not
Android APK acceptance or compatibility with other PCs.

\textbf{Handoff:} Extract the ZIP and run \code{My project.exe} in
\path{CampusEmergency} with Quest Link active and the data folder and player
DLLs beside it.

\subsection{Functional verification summary}
C\# compilation and \textbf{147 isolated Unity regression checks passed} for
the current source, including XR teleport filtering, barrier collision,
ordered recovery and two clean scene reloads. These are automated regressions,
not participant trials or full-campus acceptance (Appendix A).

\begin{tabularx}{\textwidth}{@{}p{0.8cm}YY@{}}
\toprule
\textbf{Key} & \textbf{Executed or required check} & \textbf{Result / evidence} \\
\midrule
V1 & Endpoint before Continue; persistent card; task isolation. & Pass; ordinary mode can acknowledge arrival before ordered completion. \\
V2 & Diagram, Continue, panel movement and limits. & Pass; exact labels in evidence file. \\
V3 & Optional closure, barrier, valid recovery and ordered stairs. & Pass in isolated scene; authored scenario, not participant evidence. \\
V4 & Lift proximity; muted journey. & Proximity passes; muted journey not recorded. \\
V5 & Pause, clean classroom reset and exit. & Automated reset passes; PCVR test report is summarised in V6. \\
V6 & PCVR build on lab PC. & Team-reported pass; 17 October 2026, PQ604B/604B205. \\
\bottomrule
\end{tabularx}

\newpage
% BODY PAGE 4
\section{Limitations and Evaluation Priorities}
\subsection{Operational definitions: recovery and assistance}
Phase 1 feedback requested explicit criteria for successful recovery and
acceptable assistance. The following is a \textbf{provisional Phase 3 protocol},
not a measured result or client-approved safety standard. It becomes applicable
only after an unavailable-route event is implemented.

\textbf{Successful recovery:} After the closure cue appears, the participant
recognises the closure through observable behaviour or explanation, chooses
the signposted alternative, and reaches the fictional endpoint without bypass.
Pilot a 120-second active-time window from cue presentation to endpoint entry,
excluding pauses and technical interruptions, then freeze the threshold before
the main study. Record recognition, route selection and completion separately.

\textbf{Assistance:} H0 = no route-recovery hints; H1 = at most one
non-directional prompt (e.g., ``Look again at the signs''); H2 = repeated
prompts, explicit directions or facilitator-led navigation. Log wording,
source and time, including system prompts. Standard control briefing is
separate from route help. The current arrow diagram/automatic directions must
be recorded as directional assistance, not silently counted as H0.

\textbf{Judgement:} Report completed recoveries as H0 independent, H1 acceptably
assisted, or H2 beyond the proposed allowance. Report timeouts, incomplete
attempts, technical interruptions and comfort withdrawals separately, with
denominators. Always provide safety/comfort help regardless of classification.

\subsection{Evaluation priorities and evidence}
\textbf{E1 -- Cue interpretation (R2):} Observe missed signs, hesitation and
control errors; ask which cue informed each choice.
\textbf{E2 -- Recovery/rationale (R1, R3, R4):} The closure, alternative and
safety review are now implemented. Record event times and assistance; ask for
rationale before the review and record post-review explanations separately.
\textbf{E3 -- Comfort/inclusion (R5, R6):} Observe reading/reach difficulty,
pause/exit use and discomfort; participation is voluntary.

\textbf{Study status:} Group members provided iterative hands-on feedback,
summarised in Appendix A. This is retrospective formative convenience testing,
not an independent study. The number present at each iteration, individual
attribution, exact test dates, device/build and standardised outcomes were not
logged. GitHub commit dates bound the development period but are not treated as
test timestamps. No success rate, time saving or learning gain is claimed.
Future evaluation should record anonymous task outcomes, prompts and optional
comments with consent; do not retain names, account IDs or headset recordings
by default.

\subsection{Known limitations and next priorities}
\textbf{1. Establish recovery for unfamiliar users (R1/R3/R4):} The closure,
alternative and review now exist, but there is no independent evidence that
users notice the closure, choose the alternative or understand the rationale.
\textbf{2. Establish delivery reliability (R6):} Repeat PCVR verification on
additional required lab PCs or headset/runtime combinations. Test Android
standalone deployment separately if it remains a delivery target.
\textbf{3. Pilot clarity and assistance (R2/R5):} Check mute comprehension,
turning text, fixed-panel reach and sign alignment on-device. Teammate
familiarity, a single layout and absent transfer evidence limit conclusions.

\clearpage
% REFERENCES: excluded from body limit.
\fancyhead[R]{References}
\pagenumbering{roman}
\begin{thebibliography}{9}
\bibitem{phase1}
Group 7. \textit{Requirements and Prototype Proposal: Scenario D --- Campus
Emergency Preparedness}. COMP5424 Extended Reality, 20 September 2026.
Repository copy: \path{GroupProjectDocuments/Group 7 - Group Component.pdf}.

\bibitem{brief}
COMP5424 Extended Reality. \textit{Phase 2 assignment instructions:
Purpose, Minimum Prototype Expectations and Submission Components}.
Course instructions supplied by the team; publication date/version not provided.
Includes the four-page technical-note and target-platform requirements.

\bibitem{implementation}
Group 7. \textit{Campus Emergency XR prototype}, release branch
\code{final-submission-2026-10-17}, source commit \code{9f2b3ff}, 17 October 2026.
\url{\repo}.
Implementation evidence: \path{Packages/manifest.json},
\path{ProjectSettings/ProjectVersion.txt},
\path{ProjectSettings/ProjectSettings.asset},
\path{Docs/SoundImplementation.md}, and
\path{Tools/Sound/Validation/SoundVerification.cs}.
\end{thebibliography}

\clearpage
\appendix
\fancyhead[R]{Appendix: Evidence}
\section{Evidence Locations and Submission Checks}
\subsection{Source locators}
Paths are relative to the repository root. This report describes
\code{final-submission-2026-10-17} at \code{9f2b3ff}. The Windows PCVR ZIP is a
separate build artifact produced from that source revision. The Android
standalone APK was not rebuilt for this acceptance.

\begin{tabularx}{\textwidth}{@{}p{0.8cm}Y@{}}
\toprule
\textbf{Key} & \textbf{Locator and interpretation} \\
\midrule
S1 & \path{Assets/Scenes/SampleScene.unity}: Start/End NPCs, DownStair
waypoints, LiftSign and startup menu. \\
S2 & \path{Assets/Scripts/UI/ControllerGuidePanel.cs}:
\code{LateUpdate}, \code{AdvanceRoute}, \code{TryShowEndDialogue},
\code{AdjustPanelDistance}, pause/restart/exit handlers. \\
S3 & \path{Assets/Scripts/Building/NpcGuideDialogue.cs}:
\code{ShowIntroduction}, \code{ShowTaskGuide},
\code{ShowCompletionDialogue}. \\
S4 & \path{Assets/Scripts/Audio/CampusAudioDirector.cs}:
\code{IsNearLiftSign}, \code{WarnAgainstLift}, \code{UpdateEvacuation};
\path{Assets/Scripts/Audio/LiftSafetyAudio.cs}. \\
\bottomrule
\end{tabularx}

\subsection{Executed regression evidence}
The complete saved PASS record from the isolated, no-graphics Unity project is
reproduced below. The persistent test definitions are in
\path{Tools/Sound/Validation/SoundVerification.cs}; execute using
\path{Tools/Sound/Verify-Sound.ps1} with the Unity editor path.
\textbf{Result: PASS; 147 named checks.} The suite uses an isolated fixture
and does not verify the entire campus or Quest build.

\begingroup
\small
\begin{itemize}[leftmargin=*,nosep]
\item 22 audio assets import successfully
\item NPC voice matches scene text
\item Dynamic task-list markers match baked NPC speech
\item All eight lift interaction panels use the actual sign positions on four floors
\item Lift panels and colliders are hidden when the player is away
\item Elevator NPC and End NPC locations do not count as lift sign proximity
\item Actual lift marker is detected: LiftSign 0
\item Actual lift marker is detected: LiftSign 1
\item Actual lift marker is detected: LiftSign 2
\item Actual lift marker is detected: LiftSign 3
\item Actual lift marker is detected: LiftSign 4
\item Actual lift marker is detected: LiftSign 5
\item Actual lift marker is detected: LiftSign 6
\item Actual lift marker is detected: LiftSign 7
\item Lift detection rejects other floors and positions outside 2.5 meters
\item Offset anchor resolves to actual destination
\item Training route uses corrected anchor position
\item Evacuation diagram imports as Texture2D
\item Start NPC welcome precedes the optional controller tutorial
\item Task updates do not overwrite welcome before Continue
\item Continue is an enabled XR-raycastable button
\item Continue reveals controller tutorial and dismisses welcome button
\item Task card includes the correct evacuation diagram with preserved aspect ratio
\item Task list and teleport-effect instructions appear together
\item Repeated Continue and bubble styling preserve the task card
\item Guide Nearer physically moves its panel 0.2 meters without changing height or orientation
\item Guide Farther physically reverses the Nearer adjustment
\item Guide Nearer stops at 0.65 meters and disables its button
\item Guide Farther stops at 2.2 meters and disables its button
\item Guide remains world-fixed after head rotation and player translation
\item Completing all five controller exercises opens the navigation panel
\item Post-tutorial navigation panel does not follow head rotation, walking or teleport displacement
\item Nearer adjusts a distant fixed HUD by one step instead of snapping it to the player
\item HUD Nearer physically moves its panel while preserving height and orientation
\item HUD Farther physically reverses the Nearer adjustment
\item HUD Nearer clamps to 0.65 meters
\item HUD Farther clamps to 2.2 meters
\item Review explicitly reopens the fixed guide in front of the current viewer
\item Skipping the guide also opens a world-fixed navigation panel
\item Endpoint dialogue does not trigger outside actual anchor radius
\item Endpoint dialogue does not trigger from another floor
\item Inactive endpoint NPC does not consume arrival trigger
\item Direct endpoint arrival shows configured End NPC text without prior route checkpoints
\item Endpoint dialogue leaves ordered task progress unchanged
\item Repeated endpoint arrival preserves the persistent card without replaying dialogue
\item Completion card remains visible beyond its former seven-second timeout
\item Completion card offers an optional controller-ray safety review button
\item Completion card lower edge is raised 0.4 meters above the original panel lower edge
\item Safety review explains stairs, lifts, closure awareness and the simulated endpoint without claiming reroute success
\item Opening safety review preserves persistent display, voice cooldown and ordered route progress
\item Taller safety review keeps the same raised lower edge, clear of the NPC head
\item Raised review stays below the next floor at the End NPC's configured four-meter storey height
\item Safety review body fits its panel at the configured minimum readable font size
\item Safety review remains visible beyond the former dialogue timeout
\item Back restores the configured completion message without replaying completion voice
\item Reset after safety review clears the review page and persistent message for the next experience
\item End NPC proximity does not trigger completion from another floor
\item Approaching End NPC after reset restores the completion card rather than the previous review page
\item Actual frame update still blocks endpoint completion while paused
\item Actual frame update shows permanent completion before Continue at (48.72, 12.08, -10.44)
\item Actual frame update shows permanent completion before Continue at (55.22, 12.08, -10.44)
\item Completion before Continue plays or queues the configured completion voice
\item Completion does not bypass Continue for the tutorial or advance ordered route progress
\item Emergency starts idempotently
\item Indoor alarm fades in to configured level
\item Exact emergency PA is selected
\item Speech ducks alarm for intelligibility
\item Triggered NPC preserves currently displayed task-list text
\item Triggered NPC queues the voice for the visible text
\item XRI accepts real teleport request
\item No teleport sound on request/aim/cancel phase
\item Real XRI teleport completes and moves player
\item Arrival audio is triggered by completed teleport
\item Grounded continuous movement triggers footsteps
\item NPC proximity cannot show the lift warning even when explicitly requested
\item Lift proximity detected at matching floor
\item Nearby lift marker enables its panel and XR collider
\item Wrong lift warning is anchored at the matching lift marker, not the player or NPC
\item Moving away hides lift panels and disables raycast targets
\item Pause freezes alarm fades
\item G-floor covered corridor does not count as outdoors
\item G-floor open courtyard fades alarm without false completion
\item Assembly point triggers success once
\item Assembly retains only distant alarm
\item EndEmergency clears alarm and speech queue
\item Reroute validation is registered on the real XR Interaction Manager
\item Optional reroute starts idle with no barrier or restrictions on ordinary exploration
\item Original route sign text has viewer-correct left/right and up/down orientation from all four sides
\item Start NPC task panel retains its original safe height below the next floor
\item Start reroute button moves to the right panel edge, away from the NPC head
\item Task card content restores its original unraised layout
\item Optional exercise offers Start reroute while preserving the ordinary task diagram
\item Starting reroute establishes the original route before any closure appears
\item Repeated start cannot reset an active reroute exercise
\item Exercise permits the initial corridor but explains why direct teleport to stairs is unavailable
\item Visiting the alternative corner before the closure does not count as recovery
\item Pause prevents a route closure from triggering
\item Standing below the decision corner cannot trigger closure from another floor
\item Reaching the decision corner changes the previously available route to closed exactly once
\item Remaining at the closure does not repeat the event
\item The visible closure has a static solid collider across the corridor, without a falling Rigidbody
\item A real CharacterController continuous move is stopped by the closure instead of passing through
\item Both fixed and arbitrary destinations beyond the barrier, downstairs and across the courtyard are rejected
\item The registered selection filter rejects a real stair anchor before recovery
\item A real XR ray hitting a free-area floor beyond the barrier cannot select it through the manager
\item Valid corner-to-corridor teleports are allowed without treating their flight line as walking, while courtyard landings stay blocked
\item A player near the barrier can teleport back to safety, but not through to the other side
\item Notification regression uses an actively held real XR select input
\item An actual invalid teleport gesture reports its reason once
\item Holding the same invalid teleport gesture does not repeat a popup after the former four-second cooldown
\item Releasing the teleport gesture rearms feedback for a new attempt
\item A new invalid teleport gesture can report its reason again
\item The same real XR ray and free-area floor remain selectable on the available alternative corridor
\item The opposite same-floor corridor remains available for teleporting
\item Alternative corner opens the connecting corridor, not an immediate success
\item Rejoining stairs restores normal destinations but keeps the closed barrier non-teleportable
\item Active exercise hides the answer diagram and start button, retaining environmental-sign guidance
\item Recovery alone cannot show success before the ordered stair checkpoints are complete
\item An out-of-order endpoint after recovery still cannot complete the stair route
\item The original five stair checkpoints still complete in order after reroute recovery
\item The real final arrival frame advances the last checkpoint before checking success, without a false incomplete warning
\item Recovery plus completed stair checkpoints enables the persistent endpoint completion and review
\item The final End NPC card offers Continue and Back alongside the retained safety review
\item Continue does not move, rotate, reset progress, hide the card or replay completion voice
\item Safety review uses Back to results and hides the classroom-return controls to avoid ambiguous Back buttons
\item Returning from safety review restores both finish actions
\item Dialogue reset hides both completion action buttons
\item Disabling the exercise removes its global interaction filter
\item Scene cleanup releases singleton and listeners
\item Ordinary scene entry still displays Enter System and plays menu music
\item The reloaded classroom reconnects all four tutorial input actions
\item A fresh reroute can begin in clean attempt 0
\item Restart fixture contains an active barrier, fire audio, popup and moved prop before Back
\item Back unloads the previous experience including runtime objects, old XR rig and moved props
\item Back skips Enter System and its music, opening the fresh classroom welcome instead
\item Reload restores the classroom spawn and original prop placement
\item Reload clears closure, task checkmarks, endpoint success and prior completion UI
\item Reload resets fire/global state and time while registering exactly one new route filter
\item The reloaded classroom reconnects all four tutorial input actions
\item A fresh reroute can begin in clean attempt 1
\item Restart fixture contains an active barrier, fire audio, popup and moved prop before Back
\item Back unloads the previous experience including runtime objects, old XR rig and moved props
\item Back skips Enter System and its music, opening the fresh classroom welcome instead
\item Reload restores the classroom spawn and original prop placement
\item Reload clears closure, task checkmarks, endpoint success and prior completion UI
\item Reload resets fire/global state and time while registering exactly one new route filter
\item Classroom skip is consumed once: the normal Restart action still opens Enter System
\end{itemize}
\endgroup

\subsection{Retrospective group-member formative test log}
Group members report hands-on iterations covering the full route, optional
reroute, lift-safety prompts, controller/UI use and the Back reset. The issue
notes below are consolidated from development feedback, not a contemporaneous
test protocol. Individual attribution, per-session headcount, exact test
dates and per-person task outcomes for earlier iterations were not retained.
The Windows PCVR acceptance is recorded under V6. The GitHub development
history spans 12--17 October 2026; this is not a claim that each iteration
occurred on those commit dates.

\begin{tabularx}{\textwidth}{@{}p{2.35cm}YY@{}}
\toprule
\textbf{Task / observation} & \textbf{Feedback recorded} & \textbf{Response and evidence} \\
\midrule
Controls and UI &
Navigation/tutorial panels were reported to follow the player or distance
controls appeared ineffective. &
Panels are world-fixed; Nearer/Farther move them in bounded 0.2\,m steps.
Automated checks cover transform stability and distance limits. \\
Lift safety and completion &
Lift warning appeared near the NPC rather than the lift signs; completion
feedback was also reported missing or misplaced. &
Lift prompts use the eight actual LiftSign markers. Completion is triggered by
the final anchor or End NPC and remains visible. Regression checks cover
marker proximity and endpoint behaviour. \\
Safety review &
The Review safety button was obscured by the NPC head. &
The shared completion/review card layout was raised; panel-bound checks pass
in the isolated regression suite. \\
Optional reroute &
The first route sign appeared mirrored and invalid-teleport notices recurred
frequently. &
Sign orientation was corrected; teleport validation now checks destination
and barrier crossing, and rate-limits feedback per held gesture. Regression
checks cover four viewing directions and repeated selection. \\
Back reset &
Returning to spawn left traces of the previous attempt. &
Back now reloads the scene and opens the classroom welcome directly. The
automated suite verifies two consecutive clean reloads. \\
Windows PCVR build &
The team reported successful launch and operation of all intended features on
PQ604B/604B205 via Quest Link/Air Link; no itemised task log was retained. &
\path{Build/CampusEmergency.zip}, source \code{9f2b3ff}, 17 October 2026;
single-machine PCVR result, not Android APK acceptance. \\
\bottomrule
\end{tabularx}

\subsection{Interpretation of evidence}
Automated regressions document implemented behaviour; the reported manual
result documents one Windows PCVR setup. Neither constitutes a controlled
learner study. No claims are made about learning gains, independent route
choice, comfort, or compatibility beyond the tested PC/headset configuration.

\subsection{Weekly Activity Record and Division of Work}
The project progressed from the Unity/XR foundation through environment and
interaction development to final integration and validation:
\begin{itemize}
\item \textbf{29 September--3 October:} Established the Unity/XR project
structure, configured Quest/OpenXR settings, assembled the initial campus
scene and began teleport, controller-turning and tutorial interactions
(\code{a1888f7}, \code{79e93e7}, \code{6c2beed}, \code{c54852e},
\code{d6cbf34}, \code{4bb475e}, \code{a962ae6}, \code{49bddb1}; PR \#9).
\item \textbf{5--9 October:} Added campus furniture, evacuation signs and
posters and imported asset packages
(\code{384d062}, \code{1aeddf4}, \code{703edf0}, \code{b9ee901}; PR \#13).
\item \textbf{12--16 October:} Refined the scene and UI and integrated
lift-blocking behaviour, evacuation audio, a startup menu and English
campus-resource descriptions (\code{2ff3ece}, \code{d12625f},
\code{e0b82b9}, \code{5693305}, \code{786c9b6}, \code{23e006f},
\code{89d2419}; PRs \#17--\#18).
\item \textbf{17 October:} Integrated the optional reroute exercise,
teleport-destination restrictions and clean Back/reset flow, then adjusted
the Start reroute control. The project records 147 isolated Unity regression
checks and a team-reported Windows PCVR test on PQ604B/604B205 via Quest
Link/Air Link (\code{416ca44}, \code{364e4a6}, \code{74234cb}; PR \#22).
\end{itemize}

\begin{tabularx}{\textwidth}{@{}p{2.4cm}YY@{}}
\toprule
\textbf{Member} & \textbf{Responsibilities} & \textbf{Project evidence} \\
\midrule
SUN Jiashi &
XR implementation and integration: campus scene, locomotion, teleport and
turning, interaction flow, reroute logic, reset behaviour and final integration. &
\code{d6cbf34}, \code{f15de48}, \code{364e4a6}, \code{416ca44},
\code{74234cb}; changes across \path{Assets/Scenes} and
\path{Assets/Scripts}. Phase 1 Development Group. \\
FENG Kedi &
Unity development and interaction/scene work, including platform integration
and the initial controller tutorial panel. &
Phase 1 Development Group; \code{feat/FKD--develop} (PR \#7) and
\code{feat/Cody-Scripts} (PR \#9, initial tutorial interaction panel). \\
SONG Zheyi &
Campus resource images, signs and posters, English descriptions, evacuation
audio and sound feedback; product/content coordination. &
\code{89d2419}, \code{786c9b6}; PRs \#13 and \#18. Phase 1 Product Group. \\
WENG Zhirui &
Emergency-related content and asset preparation/handoff; product requirements
and scope coordination. &
\code{a962ae6} (Module 5 content and asset handoff pack). Phase 1 Product Group. \\
ZHU Yu &
Testing and acceptance responsibilities, together with Unity asset-package
preparation. &
Phase 1 Test Group; \code{b9ee901} (Unity Asset Store packages and Git LFS). \\
CHU Silu &
Test planning, acceptance and usability feedback. &
Phase 1 Test Group; the Phase 2 feedback record covers UI, lift prompts,
reroute signs and reset behaviour. \\
\bottomrule
\end{tabularx}

\textbf{Prototype outcome:} The delivered journey includes a classroom welcome
and task briefing, a stair-based practice route to the 2F endpoint and an
optional authored reroute exercise at the 3F decision area. The reroute is a
bounded scenario, not dynamic route planning. The PCVR test covers one
reported setup; Android standalone acceptance and structured participant
evaluation remain future work.
\end{document}
